% Part: proof-theory
% Chapter: cut-elimination
% Section: ce-topmost

\documentclass[../../../include/open-logic-section]{subfiles}

\begin{document}

\olfileid[pt]{pt}{cut}{top}

\olsection{Eliminação de cortes superiores}

\begin{defn}
Considere !!a{proof}~$\pi$ que termina em uma regra \CutCS{},
\[
\AxiomC{}
\RightLabel{$\pi_1$}
\Deduce$\Gamma \fCenter \Delta, !A$
\AxiomC{}
\RightLabel{$\pi_2$}
\Deduce$!A, \Gamma \fCenter \Delta$
\RightLabel{\CutCS}
\BinaryInf$\Gamma \fCenter \Delta$
\DisplayProof
\]
e não contém nenhuma outra inferência \CutCS{}. A \emph{altura de corte}
de~$\pi$ é $\cheight{\pi} = \pheight{\pi_1} + \pheight{\pi_2}$. O
\emph{posto de corte} de~$\pi$ é $\cutrank{\pi} = \depth{!A}$.
\end{defn}

\begin{lem}\ollabel{lem:cut-adm-G3c}
A regra \CutCS{} é admissível em~\Log{G3c}. Em outras palavras, se $\pi$
é !!a{proof} que termina em uma única~\CutCS{} e, fora isso, usa apenas
regras de~$\Log{G3c}$, então existe !!a{proof}~$\pi'$ do mesmo sequente final
em~\Log{G3c}.
\end{lem}

\begin{proof}
A demonstração é feita por indução dupla na !!{height} e no posto do
corte. A !!{proof}~$\pi$ termina em:
\[
\AxiomC{}
\RightLabel{$\pi_1$}
\Deduce$\Gamma \fCenter \Delta, !A$
\AxiomC{}
\RightLabel{$\pi_2$}
\Deduce$!A, \Gamma \fCenter \Delta$
\RightLabel{\CutCS}
\BinaryInf$\Gamma \fCenter \Delta$
\DisplayProof
\]

Base da indução: $\cheight{\pi} = 0$ e $\cutrank{\pi} = 0$. Como
$\cheight{\pi} = \pheight{\pi_1} + \pheight{\pi_2} = 0$, tanto $\Gamma
\Sequent \Delta, !A$ quanto $!A, \Gamma \Sequent \Delta$ são axiomas.
Há dois casos: (a)~$!A$ não é principal em um deles, ou (b)~$!A$ é
principal em ambos. Nos dois casos mostraremos abaixo que $\Gamma
\Sequent \Delta$ deve ser um axioma (e, para isso, não precisamos da
hipótese de indução).

Distinguiremos agora vários casos possíveis conforme $\pi_1$ e $\pi_2$
sejam axiomas ou terminem em uma inferência, e conforme $!A$ seja ou não
principal nessa inferência. Os casos A e~B cobrem a base da indução. No
passo indutivo, supomos $\cheight{\pi} > 0$ ou $\cutrank{\pi} > 0$ (ou
ambos). Podemos usar a hipótese de indução: a última~\CutCS{} pode ser
eliminada de qualquer !!{proof} que termine em uma única~\CutCS{} e que
tenha posto de corte $= \cutrank{\pi}$ e !!{height} de corte $<
\cheight{\pi}$, ou tenha posto de corte~$< \cutrank{\pi}$. O caso em que
$\cheight{\pi} = 0$ (ambas as premissas são axiomas) é coberto pelos
casos~A e~B. Os casos em que $\cheight{\pi} > 0$ são cobertos pelos casos
C--D.

\paragraph{Caso A:}
Algum $!B$ ocorre tanto em $\Gamma$ quanto em~$\Delta$.
Então $\Gamma \Sequent \Delta$ é um axioma (se $!B$~for atômica), ou
tem !!a{proof}~$\pi'$ em \Log{G3c} sem \CutCS{} por
\olref[seq][ex]{prop:G3c-axioms}. Isso cobre o caso~(a) da base da
indução, pois, se $!A$ não é principal em $\Gamma \Sequent \Delta, !A$
ou não é principal em $!A, \Gamma \Sequent \Delta$, e esses sequentes
são axiomas, então alguma outra fórmula atômica~$!B$ deve ocorrer tanto
em $\Gamma$ quanto em~$\Delta$. Também cobre o passo indutivo quando uma
das premissas é um axioma no qual $!A$ não é principal (mesmo que a outra
premissa seja a conclusão de uma regra).

\paragraph{Caso B:} Ambas as premissas são axiomas, e $!A$ é principal
e, portanto, atômica. Considere a premissa esquerda, $\Gamma \Sequent
\Delta, !A$. Como $!A$ é principal, deve ocorrer em $\Gamma$ (caso
contrário, $\Gamma \Sequent \Delta, !A$ não seria um axioma). Do mesmo
modo, $!A$ deve ocorrer em $\Delta$; caso contrário, $!A, \Gamma
\Sequent \Delta$ não seria um axioma. Assim, $!A$ é atômica e ocorre
tanto em $\Gamma$ quanto em $\Delta$, isto é, $\Gamma \Sequent \Delta$
é um axioma. Isso cobre o caso~(b) da base da indução.

\begin{center}
Versões alternativas dos casos A e B
\end{center}

\paragraph{Caso A:} Uma das premissas do corte é um axioma da forma
$!C, \Gamma' \Sequent \Delta', !C$, com $!C$~atômica. Se $!C$ não é a
!!{formula} de corte~$!A$, então $!C$ deve ocorrer tanto em $\Gamma$
quanto em~$\Delta$. Nesse caso, $\Gamma \Sequent \Delta$ é um axioma e
podemos tomá-lo como~$\pi'$. Se $!C \ident !A$ é a !!{formula} de corte,
então ou (se a premissa esquerda é um axioma) $!A \in \Gamma$ e $\Gamma
= \Gamma', !A$, ou (se a premissa direita é um axioma) $!A \in \Delta$
e $\Delta = \Delta', !A$. No primeiro caso, $\pi_2$ é !!a{proof} de $!A,
!A, \Gamma' \Sequent \Delta$ e, pela admissibilidade de
\LeftR{\Contraction} (\olref[seq][inv]{prop:G3c-cont-adm}), existe
!!a{proof} de $\Gamma \Sequent \Delta$. No segundo caso, $\pi_2$ é
!!a{proof} de $\Gamma \Sequent \Delta', !A, !A$ e, pela admissibilidade
de \RightR{\Contraction}, obtemos !!a{proof} de $\Gamma \Sequent
\Delta$.

\paragraph{Caso B:} Uma das premissas é um axioma da forma
$\lfalse, \Gamma' \Sequent \Delta'$. Se a premissa esquerda é um axioma
desse tipo, então $\lfalse \in \Gamma$ e $\Gamma \Sequent \Delta$
também é um axioma. Se a premissa direita é um axioma desse tipo, então
ou $\lfalse \in \Gamma$ (caso em que $\Gamma \Sequent \Delta$ novamente
é um axioma), ou $!A \ident \lfalse$. Neste último caso, $\pi_1$ é
!!a{proof} de $\Gamma \Sequent \Delta, \lfalse$. Podemos transformá-la
em !!a{proof} de~$\Gamma \Sequent \Delta$ eliminando uma ocorrência de
$\lfalse$ do consequente de cada sequente em~$\pi_1$. (Exercício:
verifique isso por indução em~$\pheight{\pi_1}$.)

Suponha agora que nem o caso~A nem o caso~B se aplique. Nos casos
restantes, $\cheight{\pi} > 0$, isto é, pelo menos uma das premissas é a
conclusão de uma inferência.

\paragraph{Casos C e D: permutação de cortes.}
As possibilidades consideradas nos casos~C e~D têm um padrão comum.
Começamos com !!a{proof} que termina em~\CutCS{}, na qual uma premissa é
obtida por uma regra~$R$ em que a !!{formula} de corte~$!A$ não é
principal (alguma outra !!{formula}~$!D$ o é). Nessa !!{proof}, \CutCS{}
vem depois da regra~$R$ e, esquematicamente, tem a seguinte forma:
\[
\AxiomC{}
\RightLabel{$\pi_1$}
\Deduce$\Gamma \fCenter \Delta', !D, !A$
\AxiomC{}
\RightLabel{$\pi_3$}
\Deduce$!A, \Gamma, \Pi \fCenter \Delta', \Lambda$
\RightLabel{$R$}
\UnaryInf$!A, \Gamma \fCenter \Delta', !D$
\RightSubproofLabel{$\pi_2$}
\RightLabel{\CutCS}
\BinaryInf$\Gamma \fCenter \Delta', !D$
\DisplayProof
\]
Isso cobre apenas o caso em que $R$ é uma regra direita com uma premissa,
$A$ não é a !!{formula} principal de~$R$, e a !!{formula}~$!D$ que
\emph{é} principal ocorre no sucedente da premissa direita de~\CutCS.
Se $!D$ ocorre no antecedente (isto é, se $R$ é uma regra esquerda) ou
na premissa esquerda de~\CutCS, a forma é diferente, naturalmente. As
!!{formula}s em $\Pi$ e~$\Lambda$ representam as !!{formula}s ativas da
regra~$R$.

Invertemos essa ordem e consideramos !!a{proof} na qual $R$ vem depois
de~\CutCS{}:
\[
\AxiomC{}
\RightLabel{$\pi_1'$}
\Deduce$\Gamma, \Pi \fCenter \Delta', !D, \Lambda, !A$
\AxiomC{}
\RightLabel{$\pi_3'$}
\Deduce$!A, \Gamma, \Pi \fCenter \Delta', !D, \Lambda$
\RightLabel{\CutCS}
\BinaryInf$\Gamma, \Pi \fCenter \Delta', !D, \Lambda$
\RightLabel{$R$}
\UnaryInf$\Gamma, \fCenter \Delta', !D, !D$
\DisplayProof
\]
Usamos enfraquecimento que preserva a !!{height} para obter $\pi_1'$ e
$\pi_3'$ a partir de $\pi_1$ e $\pi_3$, respectivamente.

Como trocamos a ordem de \CutCS{} e~$R$, dizemos que estamos
``permutando um corte para cima através de~$R$''. Isso diminui a altura
da !!{proof} que termina na premissa direita de~\CutCS{} e, portanto,
diminui a !!{height} de corte. Isto é, podemos supor que as
sub-!!{proof}s que terminam na nova \CutCS{} têm !!{height} não maior que
as de $\pi_1$ e $\pi_3$; assim, a !!{height} de corte é reduzida
(eliminamos uma inferência à direita). A hipótese de indução agora se
aplica e nos diz que a inferência \CutCS{} pode ser eliminada. Por fim,
eliminamos a ocorrência adicional de~$!D$ usando a admissibilidade
de~$\RightR{\Contraction}$.

\paragraph{Caso C:}
Pelo menos uma entre $\pi_1$ e $\pi_2$ termina em uma regra~$R$ de uma
premissa na qual $!A$ não é principal. Há 18 casos ao todo, conforme
$!A$ não seja principal na premissa esquerda ou direita de~\CutCS{}, e
conforme qual das nove regras de uma premissa (\LeftR{\lnot},
\RightR{\lnot}, \RightR{\lor}, \LeftR{\land}, \RightR{\lif} e as quatro
regras de quantificadores) seja a regra~$R$. Consideramos apenas dois
exemplos: aqueles em que $!A$ não é principal na última inferência de
$\pi_2$, que termina em~$\RightR{\lif}$ ou em~\LeftR{\lexists}. Os outros
casos são análogos e ficam como exercícios.

Suponha que $\pi$ termine em:
\[
\AxiomC{}
\RightLabel{$\pi_1$}
\Deduce$\Gamma \fCenter \Delta', !B \lif !C, !A$
\AxiomC{}
\RightLabel{$\pi_3$}
\Deduce$!A, !B, \Gamma \fCenter \Delta', !C$
\RightLabel{$\RightR{\lif}$}
\UnaryInf$!A, \Gamma \fCenter \Delta', !B \lif !C$
\RightSubproofLabel{$\pi_2$}
\RightLabel{\CutCS}
\BinaryInf$\Gamma \fCenter \Delta', !B \lif !C$
\DisplayProof
\]
onde $\Delta = \Delta', !B \lif !C$. Aplicamos o enfraquecimento que
preserva a !!{height} (\olref[seq][adm]{prop:weak-G3c-adm}) a $\pi_1$
para obter !!a{proof} $\pi_1'$ de
$!B, \Gamma \fCenter \Delta', !B \lif !C, !C, !A$ (enfraquecemos com
$!B$ à esquerda e $!C$ à direita). De modo semelhante, aplicamos
\olref[seq][adm]{prop:weak-G3c-adm} a~$\pi_3$ para obter !!a{proof}
$\pi_3'$ de $!A, !B, \Gamma \fCenter \Delta', !B \lif !C, !C$
(enfraquecendo com $!B \lif !C$ à direita). A hipótese de indução se
aplica a
\[
\AxiomC{}
\RightLabel{$\pi_1'$}
\Deduce$!B, \Gamma \fCenter \Delta', !B \lif !C, !C, !A$
\AxiomC{}
\RightLabel{$\pi_3'$}
\Deduce$!A, !B, \Gamma \fCenter \Delta', !B \lif !C, !C$
\RightLabel{\CutCS}
\BinaryInf$!B, \Gamma \fCenter \Delta', !B \lif !C, !C$
\DisplayProof
\]
com !!{formula} de corte~$!A$: essa !!{proof} tem o mesmo posto de corte
e !!{height} de corte menor que~$\pi$. Isso nos dá !!a{proof}~$\pi_4$ em
\Log{G3c}, sem \CutCS{}, de~$!B, \Gamma \fCenter \Delta', !B \lif !C,
!C$. Aplique $\RightR{\lif}$ para obter !!a{proof} de $\Gamma \fCenter
\Delta', !B \lif !C, !B \lif !C$:
\[
\AxiomC{}
\RightLabel{$\pi_4$}
\Deduce$!B, \Gamma \fCenter \Delta', !B \lif !C, !C$
\RightLabel{$\RightR{\lif}$}
\UnaryInf$\Gamma \fCenter \Delta', !B \lif !C, !B \lif !C$
\DisplayProof
\]
Para obter !!a{proof}~$\pi'$ de $\Gamma \Sequent \Delta$, isto é,
$\Gamma \fCenter \Delta', !B \lif !C$, aplique
\olref[seq][inv]{prop:G3c-cont-adm} (a admissibilidade da contração).

Suponha agora que $\pi$ termine em:
\[
\AxiomC{}
\RightLabel{$\pi_1$}
\Deduce$\lexists[x][!B(x)], \Gamma' \fCenter \Delta, !A$
\AxiomC{}
\RightLabel{$\pi_3$}
\Deduce$!A, !B(c), \Gamma' \fCenter \Delta$
\RightLabel{$\LeftR{\lexists}$}
\UnaryInf$!A, \lexists[x][!B(x)], \Gamma' \fCenter \Delta$
\RightSubproofLabel{$\pi_2$}
\RightLabel{\CutCS}
\BinaryInf$\Gamma' \fCenter \Delta$
\DisplayProof
\]
Aplicamos novamente a hipótese de indução a
\[
\AxiomC{}
\RightLabel{$\pi_1[!B(c) \Sequent {}]$}
\Deduce$\lexists[x][!B(x)], !B(c), \Gamma' \fCenter \Delta, !A$
\AxiomC{}
\LeftLabel{$\pi_3[\lexists[x][!B(x)] \Sequent {}]$}
\Deduce$!A, \lexists[x][!B(x)], !B(c), \Gamma' \fCenter \Delta$
\RightLabel{\CutCS}
\BinaryInf$\lexists[x][!B(x)], !B(c), \Gamma' \fCenter \Delta$
\RightLabel{\LeftR{\lexists}}
\UnaryInf$\lexists[x][!B(x)], \lexists[x][!B(x)], \Gamma' \fCenter \Delta$
\DisplayProof
\]
Observe que a condição de variável própria é satisfeita, pois a condição
de variável própria de~\LeftR{\lexists} em~$\pi_2$ garante que $c$ não
ocorre em $\lexists[x][!B(x)]$, $\Gamma'$ nem~$\Delta$. Por fim, aplique
\olref[seq][inv]{prop:G3c-cont-adm}.


\paragraph{Caso D:}
Pelo menos uma entre $\pi_1$ e $\pi_2$ termina em uma regra~$R$ de duas
premissas na qual $!A$ não é principal. Há 6 casos. Consideramos aquele
em que $!A$ não é principal na última inferência de~$\pi_2$, que termina
em~$\RightR{\land}$; os outros casos são análogos.

Suponha que $\pi$ termine em:
\[
\AxiomC{}
\RightLabel{$\pi_1$}
\Deduce$\Gamma \fCenter \Delta', !B \land !C, !A$
\AxiomC{}
\RightLabel{$\pi_3$}
\Deduce$!A, \Gamma \fCenter \Delta', !B$
\AxiomC{}
\RightLabel{$\pi_4$}
\Deduce$!A, \Gamma \fCenter \Delta', !C$
\RightLabel{$\RightR{\land}$}
\BinaryInf$!A, \Gamma \fCenter \Delta', !B \land !C$
\RightSubproofLabel{$\pi_2$}
\RightLabel{\CutCS}
\BinaryInf$\Gamma \fCenter \Delta$
\DisplayProof
\]
A hipótese de indução se aplica às !!{proof}s
\[
\AxiomC{}
\RightLabel{$\pi_1'$}
\Deduce$\Gamma \fCenter \Delta', !B \land !C, !B, !A$
\AxiomC{}
\RightLabel{$\pi_3'$}
\Deduce$!A, \Gamma \fCenter \Delta', !B \land !C, !B$
\RightLabel{\CutCS}
\BinaryInf$\Gamma \fCenter \Delta', !B \land !C, !B$
\DisplayProof
\]
e
\[
\AxiomC{}
\RightLabel{$\pi_1''$}
\Deduce$\Gamma \fCenter \Delta, !B \land !C, !C, !A$
\AxiomC{}
\RightLabel{$\pi_4'$}
\Deduce$!A, \Gamma \fCenter \Delta', !B \land !C, !C$
\RightLabel{\CutCS}
\BinaryInf$\Gamma \fCenter \Delta', !B \land !C, !C$
\DisplayProof
\]
Aqui, as !!{proof}s de \Log{G3c} $\pi_1'$ e $\pi_1''$ são obtidas por
enfraquecimento que preserva a !!{height}
(\olref[seq][adm]{prop:weak-G3c-adm}) de~$\pi_1$ (uma vez enfraquecendo à direita
com $!B$, e outra com $!C$). As !!{proof}s $\pi_3'$ e~$\pi_4'$ são
obtidas de $\pi_3$ e~$\pi_4$, também por enfraquecimento que preserva a
!!{height}, a partir de $\pi_3$ e~$\pi_4$, respectivamente
(enfraquecendo com $!B \land !C$ à direita).

A hipótese de indução nos dá !!{proof}s~$\pi_5$ e~$\pi_6$ em \Log{G3c}
de $\Gamma \fCenter \Delta', !B \land !C, !B$ e de $\Gamma \fCenter
\Delta', !B \land !C, !C$. Combinamo-las em !!a{proof} de $\Gamma
\Sequent \Delta', !B \land !C, !B \land !C$ usando a
regra~$\RightR{\land}$:
\[
\AxiomC{}
\RightLabel{$\pi_5$}
\Deduce$\Gamma \fCenter \Delta', !B \land !C, !B$
\AxiomC{}
\RightLabel{$\pi_6$}
\Deduce$\Gamma \fCenter \Delta', !B \land !C, !C$
\RightLabel{\RightR{\land}}
\BinaryInf$\Gamma \fCenter \Delta, !B \land !C, !B \land !C$
\DisplayProof
\]
Para obter !!a{proof} de $\Gamma \Sequent \Delta$, isto é, $\Gamma
\fCenter \Delta', !B \land !C$, aplique
\olref[seq][inv]{prop:G3c-cont-adm} (a admissibilidade da contração).

\paragraph{Caso E: redução de cortes.}
O último caso diz respeito à situação em que tanto $\pi_1$ quanto
$\pi_2$ terminam em inferências nas quais $!A$ é principal. Há seis
casos, um para cada possível !!{main operator} de~$!A$.

Em cada caso, transformamos uma \CutCS{} com !!{formula} de corte~$!A$ em
um ou mais cortes sobre sub-!!{formula}s de~$!A$, a saber, sobre as
!!{formula}s ativas das inferências das quais~$!A$ é a !!{formula}
principal. Por isso, esses casos são frequentemente chamados de
``redução'' ou ``conversão'' da inferência de corte.

\begin{enumerate}
\item Se $!A \ident \lnot !B$, então $\pi$ termina em
\[
\AxiomC{}
\RightLabel{$\pi_1'$}
\Deduce$!B, \Gamma \fCenter \Delta$
\RightLabel{\RightR{\lnot}}
\UnaryInf$\Gamma \fCenter \Delta, \lnot !B$
\LeftSubproofLabel{$\pi_1$}
\AxiomC{}
\RightLabel{$\pi_2'$}
\Deduce$\Gamma \fCenter \Delta, !B$
\RightLabel{\LeftR{\lnot}}
\UnaryInf$\lnot !B, \Gamma \fCenter \Delta$
\RightSubproofLabel{$\pi_2$}
\RightLabel{\CutCS}
\BinaryInf$\Gamma \fCenter \Delta$
\DisplayProof
\]
Pela hipótese de indução, \CutCS{} pode ser eliminada de
\[
\AxiomC{}
\RightLabel{$\pi_2'$}
\Deduce$\Gamma \fCenter \Delta, !B$
\AxiomC{}
\RightLabel{$\pi_1'$}
\Deduce$!B, \Gamma \fCenter \Delta$
\RightLabel{\CutCS}
\BinaryInf$\Gamma \fCenter \Delta$
\DisplayProof
\]

\item
Se $!A \ident (!B \lor !C)$, então $\pi$ termina em
\[
\AxiomC{}
\RightLabel{$\pi_1'$}
\Deduce$\Gamma \fCenter \Delta, !B, !C$
\RightLabel{\RightR{\lor}}
\UnaryInf$\Gamma \fCenter \Delta, !B \lor !C$
\LeftSubproofLabel{$\pi_1$}
\AxiomC{}
\RightLabel{$\pi_2'$}
\Deduce$!B, \Gamma \fCenter \Delta$
\AxiomC{}
\RightLabel{$\pi_2''$}
\Deduce$!C, \Gamma \fCenter \Delta$
\RightLabel{\LeftR{\lor}}
\BinaryInf$!B \lor !C, \Gamma \fCenter \Delta$
\RightSubproofLabel{$\pi_2$}
\RightLabel{\CutCS}
\BinaryInf$\Gamma \fCenter \Delta$
\DisplayProof
\]
Considere a !!{proof} que termina em \CutCS,
\[
\AxiomC{}
\RightLabel{$\pi_1'$}
\Deduce$\Gamma \fCenter \Delta, !B, !C$
\AxiomC{}
\RightLabel{$\pi_2'''$}
\Deduce$!C, \Gamma \fCenter \Delta, !B$
\RightLabel{\CutCS}
\BinaryInf$\Gamma \fCenter \Delta, !B$
\DisplayProof
\]
onde $\pi_2'''$ é obtida de $\pi_2''$ por enfraquecimento que preserva a
!!{height}. Ela tem posto de corte $< \cutr{\pi}$, pois $\depth{!C} <
\depth{!B \lor !C}$; assim, a hipótese de indução fornece, em
\Log{G3c}, !!a{proof}~$\pi_1''$ de $\Gamma \fCenter \Delta, !B$. A !!{proof} que
termina em \CutCS,
\[
\AxiomC{}
\RightLabel{$\pi_1''$}
\Deduce$\Gamma \fCenter \Delta, !B$
\AxiomC{}
\RightLabel{$\pi_2'$}
\Deduce$!B, \Gamma \fCenter \Delta$
\RightLabel{\CutCS}
\BinaryInf$\Gamma \fCenter \Delta$
\DisplayProof
\]
tem posto de corte $= \depth{!B} < \depth{!B \lor !C}$; portanto, a
hipótese de indução se aplica novamente e fornece !!a{proof} $\pi'$ de
$\Gamma \fCenter \Delta$.

\item $!A \ident (!B \land !C)$. Exercício.

\item Se $!A \ident (!B \lif !C)$, então $\pi$ termina em
\[
\AxiomC{}
\RightLabel{$\pi_1'$}
\Deduce$!B, \Gamma \fCenter \Delta, !C$
\RightLabel{\RightR{\lif}}
\UnaryInf$\Gamma \fCenter \Delta, !B \lif !C$
\LeftSubproofLabel{$\pi_1$}
\AxiomC{}
\RightLabel{$\pi_2'$}
\Deduce$\Gamma \fCenter \Delta, !B$
\AxiomC{}
\RightLabel{$\pi_2''$}
\Deduce$!C, \Gamma \fCenter \Delta$
\RightLabel{\LeftR{\lif}}
\BinaryInf$!B \lif !C, \Gamma \fCenter \Delta$
\RightSubproofLabel{$\pi_2$}
\RightLabel{\CutCS}
\BinaryInf$\Gamma \fCenter \Delta$
\DisplayProof
\]
A !!{proof} que termina em \CutCS,
\[
\AxiomC{}
\RightLabel{$\pi_1'$}
\Deduce$!B, \Gamma \fCenter \Delta, !C$
\AxiomC{}
\RightLabel{$\pi_2''$}
\Deduce$!C, \Gamma \fCenter \Delta$
\RightLabel{\CutCS}
\BinaryInf$!B, \Gamma \fCenter \Delta$
\DisplayProof
\]
tem posto de corte menor que~$\pi$. A hipótese de indução fornece
!!a{proof}~$\pi_1''$ de $!B, \Gamma \fCenter \Delta$. Considere agora a
!!{proof} que termina em \CutCS,
\[
\AxiomC{}
\RightLabel{$\pi_2'$}
\Deduce$\Gamma \fCenter \Delta, !B$
\AxiomC{}
\RightLabel{$\pi_1''$}
\Deduce$!B, \Gamma \fCenter \Delta$
\RightLabel{\CutCS}
\BinaryInf$\Gamma \fCenter \Delta$
\DisplayProof
\]
Seu posto de corte também é menor que o posto de corte de~$\pi$; assim,
a hipótese de indução fornece, em \Log{G3c}, uma !!{proof}~$\pi'$ de
$\Gamma \fCenter \Delta$.
\item Se $!A \ident \lforall[x][!B(x)]$, então $\pi$ termina em
\[
\AxiomC{}
\RightLabel{$\pi_1'$}
\Deduce$\Gamma \fCenter \Delta, !B(c)$
\RightLabel{\RightR{\lforall}}
\UnaryInf$\Gamma \fCenter \Delta, \lforall[x][!B(x)]$
\LeftSubproofLabel{$\pi_1$}
\AxiomC{}
\RightLabel{$\pi_2'$}
\Deduce$!B(t), \lforall[x][!B(x)], \Gamma \fCenter \Delta$
\RightLabel{\LeftR{\lforall}}
\UnaryInf$\lforall[x][!B(x)], \Gamma \fCenter \Delta$
\RightSubproofLabel{$\pi_2$}
\RightLabel{\CutCS}
\BinaryInf$\Gamma \fCenter \Delta$
\DisplayProof
\]
Pela hipótese de indução, \CutCS{} pode ser eliminada de
\[
\AxiomC{}
\RightLabel{$\pi_1''$}
\Deduce$!B(t), \Gamma \fCenter \Delta, \lforall[x][!B(x)]$
\AxiomC{}
\RightLabel{$\pi_2'$}
\Deduce$!B(t), \lforall[x][!B(x)], \Gamma \fCenter \Delta$
\RightLabel{\CutCS}
\BinaryInf$!B(t), \Gamma \fCenter \Delta$
\DisplayProof
\]
onde $\pi_1''$ é obtida de $\pi_1$ (não de~$\pi_1'$!) por
enfraquecimento que preserva a !!{height}. (A !!{proof} tem o mesmo posto
de corte, mas !!{height} de corte menor.) Isso nos dá !!a{proof}
$\pi_2''$ de $!B(t), \Gamma \fCenter \Delta$.

A !!{proof}~$\pi_1'$ tem sequente final $\Gamma \Sequent \Delta,
!B(c)$, onde $c$ é a variável própria de uma inferência
$\RightR{\lforall}$. Logo, $c$ não ocorre em~$!B(x)$, $\Gamma$ nem
$\Delta$. Supomos que a !!{proof}~$\pi$ seja regular; portanto, $c$ é a
variável própria apenas da inferência $\RightR{\lforall}$ seguinte e só
ocorre acima dela, isto é, ocorre apenas em~$\pi_1'$ e não é variável
própria de nenhuma inferência em~$\pi_1'$. Como $c$ não ocorre em
$\pi_2$, não pode ocorrer em~$t$. Por
\olref[seq][qua]{lem:subst-regular}, a
!!{proof}~$\Subst{\pi_1'}{t}{c}$ é !!a{proof} de $\Gamma \Sequent
\Delta, !B(t)$. Aplicamos agora a hipótese de indução à !!{proof}
\[
\AxiomC{}
\RightLabel{$\Subst{\pi_1'}{t}{c}$}
\Deduce$\Gamma \fCenter \Delta, !B(t)$
\AxiomC{}
\RightLabel{$\pi_2''$}
\Deduce$!B(t), \Gamma \fCenter \Delta$
\RightLabel{\CutCS}
\BinaryInf$\Gamma \fCenter \Delta$
\DisplayProof
\]
(A hipótese de indução se aplica porque a !!{formula} de corte é~$!B(t)$;
assim, a !!{proof} tem posto de corte menor que~$\pi$.)
\item Deixamos o caso $!A \ident \lexists[x][!B(x)]$ como exercício.
\end{enumerate}
\end{proof}

\begin{prob}
Verifique o subcaso do caso~D na demonstração de
\olref{lem:cut-adm-G3c} em que $!A$ não é principal na última inferência
de~$\pi_1$, que termina em~$\LeftR{\lif}$. (Observação: parte deste
exercício consiste em enunciar claramente o que deve ser provado, isto
é, qual forma $\pi$ tem nesse caso.)
\end{prob}

\begin{prob}
Verifique o subcaso do caso~D na demonstração de
\olref{lem:cut-adm-G3c} em que $!A$ não é principal na última inferência
de~$\pi_1$, que termina em~$\LeftR{\lforall}$.
\end{prob}

\begin{prob}
Verifique o subcaso do caso~E na demonstração de
\olref{lem:cut-adm-G3c} em que $!A$ é principal nas últimas inferências
de $\pi_1$ e $\pi_2$, e $!A \ident (!B \land !C)$.
\end{prob}

\begin{prob}
Verifique o subcaso do caso~E na demonstração de
\olref{lem:cut-adm-G3c} em que $!A$ é principal nas últimas inferências
de $\pi_1$ e $\pi_2$, e $!A \ident \lexists[x][!B(x)]$.
\end{prob}

Observe que, no caso~E, a !!{height} de corte da !!{proof} que termina
em \CutCS{} e à qual aplicamos a hipótese de indução pode ser maior que
a !!{height} de corte da prova original. Por exemplo, no caso em que
$!A \ident (!B \lif !C)$, transformamos $\pi$ em !!a{proof} que envolve
duas inferências \CutCS{},
\[
\AxiomC{}
\RightLabel{$\pi_2'$}
\Deduce$\Gamma \fCenter \Delta, !B$
\AxiomC{}
\RightLabel{$\pi_1'$}
\Deduce$!B, \Gamma \fCenter \Delta, !C$
\AxiomC{}
\RightLabel{$\pi_2''$}
\Deduce$!C, \Gamma \fCenter \Delta$
\RightLabel{\CutCS}
\BinaryInf$!B, \Gamma \fCenter \Delta$
\RightSubproofLabel{$\pi_1''$}
\RightLabel{\CutCS}
\BinaryInf$!B, \Gamma \fCenter \Delta$
\DisplayProof
\]
A \CutCS{} superior entre $\pi_1'$ e $\pi_2''$ tem tanto posto de corte
quanto !!{height} de corte menores que os de~$\pi$. Mas a
!!{proof}~$\pi_1''$---resultado da aplicação da hipótese de indução---já
não tem !!{height} de corte menor que~$\pi$. Contudo, como a !!{formula}
de corte da \CutCS{} inferior é $!B$, seu \emph{posto} de corte é menor
que o posto de corte de~$\pi$.

\olref{lem:cut-adm-G3c} estabelece que podemos eliminar uma única
inferência \CutCS{} superior de !!a{proof}. Podemos mostrar que é possível
eliminar qualquer número de inferências \CutCS{} de !!a{proof}
aplicando-o sucessivamente às sub-!!{proof}s superiores que terminam
em~\CutCS{}.

\begin{thm}\ollabel{thm:cut-elim-G3c-topmost} Toda !!{proof} $\pi$ em
$\Log{G3c} + \CutCS$ pode ser transformada em !!a{proof} em~\Log{G3c}.
\end{thm}

\begin{proof}
  Por indução no número~$n$ de inferências \CutCS{} em~$\pi$. Se $n=0$,
  não há nada a fazer: $\pi$ já é !!a{proof} em~\Log{G3c}. Caso
  contrário, considere a sub-!!{proof}~$\pi'$ que termina em uma
  inferência \CutCS{} superior. Ela contém uma única regra \CutCS{} como
  última inferência. Aplique \olref{lem:cut-adm-G3c} para transformá-la
  em uma prova sem \CutCS{}~$\pi''$ e substitua a sub-!!{proof}~$\pi'$
  por~$\pi''$. O resultado é !!a{proof} que contém $n-1$ inferências
  \CutCS{}, e a hipótese de indução se aplica.
\end{proof}

\end{document}
